\chapter{The algebraic Novikov spectral sequence}
\label{ch:AANSS}

Consider the cofiber sequence
\[
\xymatrix@1{
S^{0,-1} \ar[r]^\tau & S^{0,0} \ar[r]^i & C\tau \ar[r]^-p & S^{1,-1},
}
\]
where $C\tau$ is the cofiber of $\tau$.
The inclusion $i$ of the bottom cell and projection $p$ to the top cell
are tools for comparing the $\C$-motivic Adams spectral sequence
for $S^{0,0}$ to the $\C$-motivic Adams spectral sequence for
$C\tau$.
In \cite{Isaksen14c}, the first author analyzed both spectral sequences simultaneously,
playing the structure of each against the other in order to obtain
more detailed information about both.
Then the structure of the homotopy of $C\tau$ was used
to reverse-engineer the structure of the classical Adams-Novikov spectral
sequence.

In this manuscript, we use $C\tau$ in a different, much more powerful
way, because we have a deeper understanding of the connection
between the homotopy of $C\tau$ and the structure of the
classical Adams-Novikov spectral sequence.
Namely, the $\C$-motivic spectrum $C\tau$ is an $E_\infty$-ring spectrum
\cite{Gheorghe18}.
\revv{
Here we are referring to the classical $E_\infty$-operad that 
parametrizes homotopy-coherent commutative multiplications.}

Moreover, 
the homotopy category of $C\tau$-modules is equivalent
to \revv{Hovey's stable derived category} of $BP_*BP$-comodules  \cite{GWX18} \cite{Krause18}.
By considering endomorphisms of unit objects, this comparison
of homotopy categories gives a structured explanation for the
identification of the homotopy of $C\tau$ and the
classical Adams-Novikov $E_2$-page.

From a computational perspective, there is an even better connection.
Namely, the algebraic Novikov spectral sequence for computing
the Adams-Novikov $E_2$-page 
\cite{Novikov67} \cite{Miller75} 
is identical to the $\C$-motivic Adams
spectral sequence for computing the homotopy of $C\tau$ \cite{GWX18}.
This rather shocking, and incredibly powerful, 
identification of spectral sequences allows us to transform
purely algebraic computations directly into information about
Adams differentials for $C\tau$.
Finally, naturality along the inclusion $i$ of the bottom cell
and along the projection $p$ to the top cell allows us to
deduce information about Adams differentials for $S^{0,0}$.

Due to the large quantity of data, we do not explicitly 
describe the structure of the Adams spectral sequence
for $C\tau$ in this manuscript.  We refer the interested
reader to the charts in \cite{IWX19}, which provide details
in a graphical form.

\section{Naming conventions}
\label{sctn:Ctau-naming}

Our naming convention for elements of the algebraic Novikov
spectral sequence (and for elements of the Adams-Novikov
spectral sequence) differs from previous approaches.
Our names are chosen to respect the inclusion $i$ of the
bottom cell and the projection $p$ to the top cell.
Specifically, if $x$ is an element of the
$\C$-motivic Adams $E_2$-page for $S^{0,0}$, then we use the same
letter $x$ to indicate its image 
$i_*(x)$ in the Adams $E_2$-page for $C\tau$.
It is certainly possible that $i_*(x)$ is zero, but
we will only use this convention in cases where 
$i_*(x)$ is non-zero, i.e., when $x$ is not a multiple of $\tau$.

On the other hand, if $x$ is an element of the $\C$-motivic
Adams $E_2$-page for $S^{0,0}$ such that $\tau x$ is zero,
then we use the symbol $\overline{x}$ to indicate an element
of $p_*^{-1} (x)$ in the Adams $E_2$-page for $C\tau$.
There is often more than one possible choice for $\overline{x}$,
and the indeterminacy in this choice equals the image of
$i_*$ in the appropriate degree.  We will not usually be explicit
about these choices.  However, potential confusion can arise
in this context.  For example, it may be the case that
one choice of $\overline{x}$ supports an $h_1$ extension,
while another choice of $\overline{x}$ supports an $h_2$ extension,
but there is no possible choice of $\overline{x}$ that simultaneously
supports both extensions.  (The authors dwell on this point
because this precise issue has generated confusion about specific
computations.)

\section{Machine computations}
\label{sctn:AANNSS-machine}

We have analyzed the algebraic Novikov spectral sequence
by computer in a large range.  Roughly speaking, our algorithm
computes a Curtis table for a minimal resolution. 
Significant effort went into optimizing the linear algebra
algorithms to complete the computation in a reasonable amount of time.
The data is available at \cite{IWX19}.
See \cite{Wang20} for a discussion of the implementation.

Our machine computations give us a full description of the
additive structure of the algebraic Novikov $E_2$-page, together
with all $d_r$ differentials for $r \geq 2$.  
It thus yields a full description of the additive structure
of the algebraic Novikov $E_\infty$-page.

Moreover, the data also gives full information about
multiplication by $2$, $h_1$, and $h_2$ in the 
Adams-Novikov $E_2$-page for the classical sphere spectrum,
which we denote by $H^*(S; BP)$.

We have also conducted machine computations of the
Adams-Novikov $E_2$-page for the classical cofiber of $2$, which
we denote by $H^*(S/2; BP)$.
Note that $H^*(S; BP)$ is the homology of a differential
graded algebra (i.e., the cobar complex) that is free as a $\Z_2$-module.
Therefore, $H^*(S/2; BP)$ 
is the homology of this differential graded algebra modulo $2$.
We have computed this homology by machine, including full information
about multiplication by $h_1$, $h_2$, and $h_3$.
These computations are related by a long exact sequence 
\[
\xymatrix@1{
\cdots \ar[r] & H^*(S; BP) \ar[r]^j & H^*(S/2; BP) \ar[r]^q &
H^*(S; BP) \ar[r] & \cdots 
}
\]


Because $h_2^2$, $h_3^2$, $h_4^2$, and $h_5^2$ are annihilated by
$2$ in $H^*(S; BP)$, there are classes
$\widetilde{h_2^2}$, $\widetilde{h_3^2}$, $\widetilde{h_4^2}$, and 
$\widetilde{h_5^2}$ in $H^*(S/2; BP)$ such that
$q(\widetilde{h_i^2})$ equals $h_i^2$ for $2 \leq i \leq 5$.
We also have full information about multiplication by
$\widetilde{h_2^2}$, $\widetilde{h_3^2}$, $\widetilde{h_4^2}$, and 
$\widetilde{h_5^2}$ in $H^*(S/2; BP)$

This multiplicative information allows use to determine some of the
Massey product structure in the Adams-Novikov $E_2$-page for the
sphere spectrum.  There are several cases to consider.

First, let $x$ and $y$ be elements of $H^*(S;BP)$.
If the product $j(x) j(y)$ is non-zero in $H^*(S/2;BP)$, then 
$x y$ must also be non-zero in $H^*(S; BP)$.

In the second case, let $x$ be an element of $H^*(S; BP)$, and let
$\widetilde{y}$ be an element of $H^*(S/2; BP)$ such that
$q(\widetilde{y}) = y$.
If the product $x \cdot \widetilde{y}$ is non-zero in
$H^*(S/2; BP)$ and equals $j(z)$ for some $z$ in
$H^*(S; BP)$, then $z$ belongs to the Massey product
$\langle 2, y, x \rangle$.  This follows immediately from the
relationship between Massey products and the multiplicative
structure of a cofiber, as discussed in \cite{Isaksen14c}*{Section 3.1.1}.

Third, let $\widetilde{x}$ and $\widetilde{y}$ be elements of
$H^*(S/2; BP)$ such that 
$q(\widetilde{x}) = x$, $q(\widetilde{y}) = y$, and
$q(\widetilde{x} \cdot \widetilde{y}) = z$.
Then $z$ belongs to the Massey product
$\langle x, 2, y \rangle$ in $H^*(S; BP)$.  This follows
immediately from the multiplicative snake lemma \ref{lem:mult-snake}.

\begin{ex}
Computer data shows that 
the product $\widetilde{h_4^2} \cdot \widetilde{h_5^2}$ does not equal
zero in $H^*(S/2;BP)$.  This implies that the
Massey product $\langle h_4^2, 2, h_5^2 \rangle$ does not contain zero in
$H^*(S; BP)$, which in turn implies that
the Toda bracket $\langle \theta_4, 2, \theta_5 \rangle$
does not contain zero in $\pi_{93,48}$.
\end{ex}

\begin{remark}
let $\widetilde{x}$ and $\widetilde{y}$ be elements of
$H^*(S/2; BP)$ such that 
$q(\widetilde{x})$ and $q(\widetilde{y})$ equal $x$ and $y$, and
such that
$\widetilde{x} \cdot \widetilde{y}$
equals $j(z)$ for some $z$ in $H^*(S; BP)$.
It appears that $z$ has some relationship to the fourfold Massey
product $\langle 2, x, 2, y \rangle$, but we have not made this precise.
\end{remark}

\begin{lemma}[Multiplicative Snake Lemma]
\label{lem:mult-snake}
Let $A$ be a differential graded algebra that has no $2$-torsion, 
and let $H(A)$ be its homology.  Also let $H(A/2)$ be the homology
of $A/2$, and let $\delta: H(A/2) \map H(A)$ be the boundary
map associated to the short exact sequence
\[
\xymatrix@1{
0 \ar[r] & A \ar[r]^2 & A \ar[r] & A/2 \ar[r] & 0.
}
\]
Suppose that $a$ and $b$ are elements of $H(A/2)$ such that 
$2 \delta(a) = 0$ and $2 \delta(b) = 0$ in $H(A)$.
Then the Massey product $\langle \delta(a), 2, \delta(b) \rangle$
in $H(A)$
contains $\delta(a b)$.
\end{lemma}

\begin{proof}
We carry out a diagram chase in the spirit of the snake lemma.
Write $\partial$ for the boundary operators in $A$ and $A/2$.

Let $x$ and $y$ be cycles in $A/2$ that represent $a$ and $b$
respectively.
Let $x'$ and $y'$ be elements in $A$ that reduce to $x$ and $y$.
Then $\partial x'$ and $\partial y'$ reduce to zero in $A/2$
because $x$ and $y$ are cycles.
Therefore, $\partial x' = 2 \widetilde{x}$ and
$\partial y' = 2 \widetilde{y}$ 
for some $\widetilde{x}$ and $\widetilde{y}$ in $A$.

By definition of the boundary map,
$\delta(a)$ and $\delta(b)$ are represented by $\widetilde{x}$
and $\widetilde{y}$.  By the definition of Massey products,
the cycle
$\widetilde{x} y' + x' \widetilde{y}$ is contained in 
$\langle \delta(a), 2, \delta(b) \rangle$.

Now we compute $\delta(a b)$.
Note that $x' y'$ is an element of $A$ that reduces to $a b$.
Then 
\[
\partial (x' y') = \partial(x') y' + x' \partial(y') =
2 (\widetilde{x} y' + x' \widetilde{y}).
\]
This shows that $\delta(a b)$ is represented by
$\widetilde{x} y' + x' \widetilde{y}$.
\end{proof}


\section{$h_1$-Bockstein spectral sequence}
\label{sctn:h1-Bockstein}

The charts in \cite{IWX19} show graphically 
the algebraic Novikov spectral sequence,
i.e., the Adams spectral sequence for $C\tau$.  Essentially all of
the information in the charts can be read off from machine-generated data.
This includes hidden extensions in the $E_\infty$-page.

One aspect of these charts requires further explanation.
The $\C$-motivic Adams $E_2$-page for $C\tau$ contains a large
number of $h_1$-periodic elements, i.e.,
elements that support infinitely many $h_1$ multiplications.
The behavior of these elements is entirely understood
\cite{GI15}, at least up to many multiplications by $h_1$,
i.e., in an $h_1$-periodic sense.

On the other hand, it takes some work to \revv{``delocalize"}
this information.  For example, we can immediately deduce from
\cite{GI15} that $d_2(h_1^k e_0) = h_1^{k+2} d_0$ for large
values of $k$, but that does not necessarily determine
the behavior of Adams differentials for small values of $k$.

The behavior of these elements is a bit subtle in another sense, 
as illustrated
by Example \ref{ex:_c0e0}.

\begin{ex}
\label{ex:_c0e0}
Consider the $h_1$-periodic element $\ol{c_0 e_0}$
\revv{in the algebraic Novikov spectral sequence}.
Machine computations tell us that this element supports a 
$d_2$ differential, but there is more than one possible value
for $d_2(\ol{c_0 e_0})$ because of the presence of
both $h_1^2 \ol{c_0 d_0}$ and $P e_0$.

In fact, $d_2(\ol{c_0 d_0})$ equals $P d_0$, and
$d_2(P e_0)$ equals $P h_1^2 d_0$.  Therefore,
$P e_0 + h_1^2 \ol{c_0 d_0}$ is the only non-zero $d_2$ cycle,
and it follows that 
$d_2(\ol{c_0 e_0})$ must equal $P e_0 + h_1^2 \ol{c_0 d_0}$.

\revv{
The careful reader will note that $d_2(\ol{c_0 e_0})$ is
not shown on the algebraic Novikov chart in \cite{IWX19}.
As discussed in \cite{IWX19}*{Section 4}, the $h_1$-periodic
differentials are not shown for legibility.
Instead, the differential is shown in the $h_1$-Bockstein
spectral sequence chart of \cite{IWX19}, up to higher
powers of $h_1$.
}
\end{ex}


In higher stems, it becomes more and more difficult to determine
the exact values of the Adams $d_2$ differentials
on $h_1$-periodic classes.  Eventually, these complications
become unmanageable because they involve sums of many monomials.

Fortunately, we only need concern ourselves with the Adams
$d_2$ differential in this context.
The $h_1$-periodic $E_3$-page equals the $h_1$-periodic
$E_\infty$-page, and the only non-zero classes are well-understood
$v_1$-periodic families running along the top of the Adams chart.

Our solution to this problem, as usual, is to introduce a filtration
that hides the \revv{higher order terms}.
In this case, we filter by powers of $h_1$.  The effect is that
terms involving higher powers of $h_1$ are ignored, and the
formulas become much more manageable.

This $h_1$-Bockstein spectral sequence starts with an $E_0$-page,
because there are some differentials that do not increase
$h_1$ divisibility.  For example, we have
Bockstein differentials $d_0(\ol{h_1^2 e_0}) = \ol{h_1^4 d_0}$
and $d_0(\ol{c_0 d_0}) = P d_0$, reflecting the Adams
differentials $d_2(\ol{h_1^2 e_0}) = \ol{h_1^4 d_0}$
and $d_2(\ol{c_0 d_0}) = P d_0$.

There are also plenty of higher $h_1$-Bockstein differentials,
such as $d_2(e_0) = h_1^2 d_0$, and
$d_7(e_0^2 g) = M h_1^8$.

\begin{remark}
Beware that filtering by powers of $h_1$ changes the multiplicative
structure in perhaps unexpected ways.
For example, $P h_1$ and $d_0$ are not $h_1$-multiples, so their
$h_1$-Bockstein filtration is zero.  One might expect their
product to be $P h_1 d_0$, but the 
$h_1$-Bockstein filtration of this element is $1$.
Therefore, $P h_1 \cdot d_0$ equals $0$ in the
$h_1$-Bockstein spectral sequence.

But not all $P h_1$ multiplications are trivial in the
$h_1$-Bockstein spectral sequence.  For example, we have
$P h_1 \cdot \ol{c_0 d_0} = \ol{P h_1 c_0 d_0}$ because
the $h_1$-Bockstein filtrations of all three elements are zero.
\end{remark}

In Example \ref{ex:_c0e0}, we explained that there is an
Adams differential $d_2(\ol{c_0 e_0}) = P e_0 + h_1^2 \ol{c_0 d_0}$.
When we throw out higher powers of $h_1$, we obtain
the $h_1$-Bockstein differential $d_0(\ol{c_0 e_0}) = P e_0$.
We also have an $h_1$-Bockstein differential
$d_0(\ol{c_0 d_0}) = P d_0$.

The first four charts in \cite{IWX19} show graphically how this 
$h_1$-Bockstein spectral sequence plays out in practice.
The main point is that the $h_1$-Bockstein
$E_\infty$-page reveals which (formerly) $h_1$-periodic classes
contribute to the Adams $E_3$-page for $C\tau$.
